% =============================================================================
% Section 4: The structural reduction
% Positive expository story; clean derivation of 96 → 41 + 143-collapse.
% =============================================================================

\section{The structural reduction}
\label{sec:structural-reduction}

This section establishes the positive content of the paper: the
sequence of structural reductions that take the original
$\tauf$-barrier search over $\N$ and reduce it to 41 residue classes
modulo $323 = 17 \cdot 19$, with all reductions Lean-verified at the
level of the Mathlib kernel. The picture that emerges is that any
counterexample to \cref{prob:e647}, if one exists, is forced into a
small, highly structured slice of $\Z$.

\subsection{The Stage-1 sieve and the 96 survivors}
\label{sec:stage1-sieve}

Following the program of \cite{ErdosGraham1980} and a Lean
formalization developed in this project (Lean source
\texttt{Erdos647ResiduePartitionStage1.lean}, available with the paper's
supplementary materials), the Stage-1 sieve eliminates a residue
class $r \pmod M$ from
contention as a potential counterexample whenever there exist
$d \in \mathcal{D}_{12}$ and $q \in \{11, 13, 17, 19\}$ with
$d r \equiv 1 \pmod q$, where
\[
\mathcal{D}_{12} \;=\; \{105, 140, 210, 252, 280, 315, 420, 504, 630,
840, 1260, 2520\}
\]
is the canonical 12-element coefficient set. The condition encodes
the existence of an early shift $k$ with $\tauf(2520(Ms + r) - k) > k+2$
forced by 11-, 13-, 17-, or 19-divisibility.

\begin{theorem}[Stage-1 sieve, restated; Lean-verified]
\label{thm:stage1-restated}
Of the $46189$ residue classes mod $M$, exactly $96$ \emph{survive}
the Stage-1 sieve --- i.e., admit no $(d, q) \in \mathcal{D}_{12}
\times \{11, 13, 17, 19\}$ with $d r \equiv 1 \pmod q$. Any
counterexample $n > 24$ to \cref{prob:e647} satisfies $n \equiv r
\pmod M$ for one of these 96 classes.
\end{theorem}

The Lean formalization in
\texttt{lean/Erdos647ResiduePartitionStage1.lean} encodes the sieve
as a decidable predicate on $\Z / M\Z$, with the survivor set
checked by \texttt{native\_decide}. The proof that any
counterexample lies in this set is the chain
\texttt{Bridge2.IsErdos647} $\Rightarrow$ \texttt{survives}
established in \texttt{lean/Bridge2.lean} (theorem
\texttt{bridge\_isErdos647\_to\_sieve}).

\subsection{The universal free-prime lift}
\label{sec:v23-universal}

A subsequent reduction observes that all 96 Stage-1 survivors share
a common $v_{23}$-rich sub-AP that is itself eliminated by a
Hensel-style lift at the prime $23$.

\begin{theorem}[Universal $v_{23}$ partial closure, restated;
Lean-verified]
\label{thm:v23-restated}
For each $r$ among the 96 Stage-1 survivors, there exists an
explicit residue $r' \pmod{Md}$ ($d = 1058 = 2 \cdot 23^2$) such that
no $n \equiv r' \pmod{Md}$ satisfies the $\tauf$-barrier inequality;
moreover, the elimination is uniform in the sense that the
$v_{23}$-rich sub-AP at every residue $r$ is closed by the same
elementary divisor count $\tauf(d \cdot p) > 4$ for any prime
$p \notin \{2, 23\}$.
\end{theorem}

The Lean proof is generated by the Hensel rootline tactic of
\texttt{lean/RequestProject/HenselRootline.lean}, instantiated 96
times in \texttt{lean/generated/V23All96.lean}. The universality is
the key feature: the $v_{23}$ lift does not depend on the
post-Stage-1 partition into 55 Lean-closed and 41 open classes ---
it acts on the entire 96.

\begin{remark}[The $55/41$ split is not a structural distinction]
\label{rem:5541-not-structural}
A historical residue partition divides the 96 Stage-1 survivors into
55 ``Lean-closed'' classes (eliminated by various structural
arguments through Bridge3 and the rootline framework) and 41
``open'' classes that resist further closure within the
finite-local toolkit. We caution that this split is best understood
as an artifact of which closure certificates have been produced to
date, rather than as a structural distinction inherent to the
problem. The universal $v_{23}$ partial closure
(\cref{thm:v23-restated}) acts uniformly across all 96 without
distinguishing the 55 from the 41; subsequent reductions
(\cref{thm:143-collapse}) likewise apply uniformly. The
characterization developed in this paper takes the residual 41 as
the operational frontier.
\end{remark}

\subsection{The 143-collapse}
\label{sec:143-collapse}

The third and most striking reduction concerns divisibility by the
two smallest primes of $M$.

\begin{theorem}[143-collapse, restated; Lean-verified]
\label{thm:143-collapse-restated}
Every $r \in \Ropen$ satisfies $11 \mid r$ and $13 \mid r$, hence
$143 = 11 \cdot 13$ divides $r$.
\end{theorem}

\begin{proof}[Proof in Lean]
The set $\Ropen$ is the explicit 41-element subset of $\Z / M\Z$
listed in \cref{thm:143-collapse}. The claim
$\forall r \in \Ropen,\, 143 \mid r$ is decidable; the Lean theorem
\texttt{openResidues41\_143\_collapse} (in
\texttt{lean/Erdos647\_143Collapse.lean}) proves it via
\texttt{native\_decide}. The Mathlib decidability of integer
divisibility makes the entire proof a single tactic invocation.
\qed
\end{proof}

\begin{remark}[A first-principles derivation]
\label{rem:143-collapse-first-principles}
The 143-collapse is not an accident. For each $q \in \{11, 13\}$,
$\gcd(2520, q) = 1$, so for every $r \not\equiv 0 \pmod q$ there is a
unique $k_r \in \{1, \ldots, q-1\}$ with $2520 r - k_r \equiv 0 \pmod q$.
The Stage-1 sieve uses these forced divisibilities to eliminate every
$r$ with $r \not\equiv 0 \pmod q$. Hence only residues with
$r \equiv 0 \pmod{11 \cdot 13} = \pmod{143}$ survive. The same
argument does not apply at $q \in \{17, 19\}$ because the
12-coefficient set $\mathcal{D}_{12}$ does not exhaust $\Z / 17\Z$
or $\Z / 19\Z$. The observation appears as a free corollary of the
Stage-1 sieve once one looks for it.
\end{remark}

\subsection{The normalized $q = 323$ frontier}
\label{sec:q323-normalization}

Combining the 143-collapse with the canonical Chinese Remainder
decomposition
\[
M \;=\; 11 \cdot 13 \cdot 17 \cdot 19 \;=\; 143 \cdot 323
\]
yields the working normalization of the rest of the paper. Every
$r \in \Ropen$ has the form $r = 143 \cdot u_r$ for a unique
$u_r \in \Z / 323\Z$; the set
\[
\Uset \;=\; \bigl\{ u_r \,:\, r \in \Ropen \bigr\}
\;\subset\; \Z / 323\Z
\]
has $|\Uset| = 41$, and is given explicitly by
\[
\Uset = \{0, 6, 9, 12, 17, 18, 34, 37, 43, 57, 63, 69, 74, 85, 88,
91, 94, 114, 119, 120, 131, 139, 142, 145, 148,
\]
\[
170, 171, 176, 190, 196, 199, 205, 207, 221, 224, 227, 247, 256,
264, 272, 309\} .
\]
Equivalently, any potential counterexample $n$ to
\cref{prob:e647} has the form
\[
n \;=\; 143 \cdot Y , \qquad Y \in \N , \qquad Y \bmod 323 \in \Uset.
\]
The shift form, on which the rest of the analysis operates, is
\[
2520 \cdot n - k \;=\; 360360 \cdot Y - k , \qquad
360360 \;=\; \lcm(1, 2, \ldots, 15) .
\]

\begin{corollary}[Full divisibility of the candidate]
\label{cor:full-divisibility}
Let $n_{\rm act} > 84$ be a counterexample to \cref{prob:e647}.
Combining the Stage-0 divisibility $2520 \mid n_{\rm act}$ (proved in
the companion main preprint as theorem~(R1), Lean: \texttt{erdos647\_div\_2520})
with the 143-collapse \cref{thm:143-collapse-restated}, we obtain
\[
360360 \;=\; 2520 \cdot 143 \;\big|\; n_{\rm act} .
\]
Writing $n_{\rm act} = 2520 \cdot m$, the search variable $m$ satisfies
\[
m \bmod M \;\in\; 143 \cdot \Uset \pmod{M} ,
\qquad M = 46189 ,
\]
where $143 \cdot \Uset \pmod{M} = \{0, 858, 1287, 1716, 2431, \ldots, 44187\}$
is the explicit 41-element set of residues lifted from $\Uset \subset \Z/323\Z$.
\end{corollary}

\subsection{Structural consequences of the normalization}
\label{sec:q323-consequences}

The normalization to $q = 323$ has several immediate consequences
that make the rest of the analysis tractable.

\begin{itemize}[leftmargin=*]
\item \textbf{Forced low-prime structure.} For every $1 \le k \le 15$
  with $k \mid 360360$, one has $360360 Y - k = k \cdot G_k(Y)$
  with $G_k(Y) = (360360 / k) Y - 1$. This factorization is exact
  and pulls a known divisor of $k$ out of the shift expression. For
  $k = 2, 4, 6, 8, 12$ the cofactor $G_k(Y)$ must be prime to
  satisfy the protected window $\tauf$-budget, since $\tauf(k) \cdot
  \tauf(\text{cofactor}) \le k + 2$ and $\tauf(k) = 2, 3, 4, 4, 6$
  respectively leaves cofactor budget at most $\lceil (k+2)/\tauf(k)
  \rceil = 2, 2, 2, 2, 2$, i.e.\ strict primality. The
  protected-window constraint is therefore strikingly rigid for
  small $k$ (cf.\ \cref{sec:csp-audit}).

\item \textbf{17- and 19-projection structure.} Reducing $\Uset$
  modulo $17$ and modulo $19$ yields, respectively,
  $\Uset \bmod 17 = \{0, 1, 3, 6, 9, 12\}$ (size 6, missing 11
  classes) and $\Uset \bmod 19 = \{0, 5, 6, 9, 12, 15, 17, 18\}$
  (size 8, missing 11 classes). The product of the projection sizes
  is $6 \cdot 8 = 48$, but $|\Uset| = 41$: the joint distribution
  is one-per-cell, so $\Uset$ misses 7 specific cells of the
  $\mathrm{proj}_{17} \times \mathrm{proj}_{19}$ product, namely
  $\{188, 233, 275, 278, 281, 284, 290\}$. The classification of
  these 7 missing cells is identified as Open Problem~5 in
  \cref{sec:open}.

\item \textbf{Search-space reduction.} The original problem ranges
  over $\N$; the structural reduction restricts the search to
  $41 / 323$ of $\N$, a factor of $\approx 7.88$. Combined with the
  $1 / 143$ density factor from the 143-collapse, the effective
  search density is $41 / (143 \cdot 323) = 41 / 46189 \approx
  8.88 \times 10^{-4}$ relative to all of $\N$. Empirical
  verification programs (cf.\ \cref{sec:empirical-detail}) operate
  in this 41-class slice rather than over the full residue
  spectrum.
\end{itemize}

\subsection{Summary of the reduction chain}
\label{sec:reduction-summary}

The complete reduction chain from $\N$ to the operational
frontier is:
\[
\N \xrightarrow{\text{Stage-1 sieve}} 96 \text{ residues mod } M
\xrightarrow{\text{$v_{23}$ partial}} 96 \text{ classes (sub-AP eliminated)}
\xrightarrow{\text{143-collapse}} \Uset \subset \Z / 323\Z, \;\; |\Uset| = 41 .
\]
All three reductions are independently Lean-verified at the level
of the Mathlib kernel; the certificates are
\texttt{Bridge2.survives} (the residue-survives Boolean predicate
verified by \texttt{bridge\_isErdos647\_to\_sieve}),
\texttt{v23\_closure\_r*} (96 instances in
\texttt{lean/generated/V23All96.lean}, each invoking
\texttt{Erdos647V23Hensel.hensel\_rootline\_closure}), and
\texttt{Erdos647\_143Collapse.openResidues41\_143\_collapse},
respectively. The remaining work --- the finite-local audit
(\cref{sec:csp-audit}), the analytic obstruction analysis
(\cref{sec:obstructions}), and the empirical verification
(\cref{sec:empirical-detail}) --- operates on the 41-class
normalized frontier.

% =============================================================================
% End of section 4
% =============================================================================
